\chapter{Experimental Setup}
\label{Setup:chap:chap6}
\graphicspath{{15.Chapter6/Images/}}

This chapter pins down the empirical setting in which FIP-FL is run: which datasets we use, how the data is split among the clients, what attack the adversary mounts, which numbers we look at, and the implementation and hardware stack that the experiments actually live on.

\section{Datasets}
\label{sec:setup-datasets}

The evaluation grid runs on two benchmarks --- one image task and one intrusion-detection task. The high-level statistics are listed in Table~\ref{tab:datasets}.

\begin{table}[H]
\centering
\caption{Datasets used in the evaluation of FIP-FL.}
\label{tab:datasets}
\begin{tabular}{@{}lcccc@{}}
\toprule
\textbf{Dataset} & \textbf{Modality} & \textbf{Classes} & \textbf{Features} & \textbf{Samples} \\
\midrule
MNIST           & Image   & 10 & $784$ & $60{,}000 / 10{,}000$ \\
KDDCup99        & Tabular & 5  & $41$  & $\approx 494$K ($10\%$) \\
\bottomrule
\end{tabular}
\end{table}

\noindent\textbf{MNIST.} The MNIST corpus has $60{,}000$ training and $10{,}000$ test grayscale images at $28\times28$, labelled over the digits $\{0,1,\ldots,9\}$. Every image is flattened to a $784$-dimensional vector and rescaled into the unit interval. The same dataset is the headline benchmark in Yazdinejad et al.~\cite{yazdinejad2024robust}, which makes it the direct comparison point against the auditor-based baseline.

\noindent\textbf{KDDCup99.} KDDCup99 is the canonical network-intrusion benchmark. The full archive contains around $4.9$ million labelled connection records; the experiments here use the standard $10\%$ subset, which is roughly $494{,}000$ records. Each record exposes $41$ features mixing traffic counters and protocol flags, and the labels are collapsed into five classes --- \textsc{Normal}, \textsc{DoS}, \textsc{Probe}, \textsc{R2L} and \textsc{U2R}. Categorical fields are one-hot encoded and every feature is standardised. The class distribution is strongly imbalanced, which is exactly what makes this dataset a useful stress test: it forces the audit to behave even when the priors are far from uniform.

\noindent\textbf{Model Architecture.} For both datasets, a single-layer softmax logistic regression classifier is used. This choice is consistent with the model used in the SenSys 2026 poster as well as the auditor-based baseline. The total number of parameters is $d=(\text{input dim})\cdot C+C$, where the final term corresponds to the bias.

A compact model is selected for two main reasons. First, Paillier encryption with $\lambda=1024$ bits is still practical when the parameter dimension remains in the range of a few thousand. Second, the convexity and $L$-smoothness assumptions required in Theorem~\ref{thm:fip-convergence} are naturally satisfied by softmax logistic regression. In contrast, for deep neural networks these assumptions would hold only approximately.

\section{Data Partitioning}
\label{sec:setup-partition}

Each dataset is sliced into $N=10$ disjoint client shards, and I run the experiments under two partition regimes. In the IID regime the training set is shuffled once and dealt out into equal-sized shards, so each client ends up with roughly $|\mathcal{D}_{\text{train}}|/N$ samples and the expected label histogram on every shard matches the global one.

For the non-IID regime I use label skew. With a $C$-class task, client $i$ gets a primary label set $\mathcal{L}_{i}\subset\{0,1,\ldots,C-1\}$ of size $\lceil C/2\rceil$, and the overlap between successive clients is $\lfloor C/4\rfloor$. The client data are assigned using an $80/20$ rule. For each client, $80\%$ of the data comes from its primary labels, while the remaining $20\%$ is sampled from the full training set. This partitioning strategy ensures that each client still retains a small number of samples from every class. The setup follows the evaluation protocol used in SenSys 2026 and is also similar to the non-IID configuration considered by Yazdinejad et al.~\cite{yazdinejad2024robust}.

At the same time, the partition is deliberately challenging because even honest clients may produce different descent directions. As a result, the audit mechanism must tolerate normal data heterogeneity while still being able to identify malicious updates reliably.

\section{Attack Models}
\label{sec:setup-attacks}

Two pre-training label-flipping attacks make up the empirical sweep. Neither needs the adversary to know anything about Paillier or about the FIP-FL audit logic; both are plain data-side manipulations.

\noindent\textbf{Untargeted label flipping.} A malicious client relabels each local sample by a fixed cyclic offset,
\begin{equation}
    y_{k}^{(i,\,\text{mal})} \;=\; (y_{k}^{(i)} + \kappa) \bmod C .
\end{equation}
On MNIST I take $\kappa=5$, which gives a no-fixed-point permutation of the ten digits, and on KDDCup99 I take $\kappa=2$ because the task only has five classes. The aim of this attack is to drag overall accuracy down by pulling each malicious client's local gradient toward a deliberately wrong classification objective.

\noindent\textbf{Targeted label flipping.} A malicious client touches only a single source class $c_s$ and remaps it to a chosen target class $c_t$,
\begin{equation}
    y_{k}^{(i,\,\text{mal})} \;=\; \begin{cases}
        c_{t} & \text{if } y_{k}^{(i)} = c_{s}, \\
        y_{k}^{(i)} & \text{otherwise.}
    \end{cases}
\end{equation}
On MNIST the chosen source--target pair is $(c_s,c_t)=(0,7)$, and on KDDCup99 I use the analogous pair drawn from the two classes that dominate in the skew regime. The point of this attack is to break one semantic class while leaving overall accuracy looking reasonable, so the relevant diagnostic metric is the target-class accuracy rather than the overall one.

\noindent\textbf{Malicious-client rates.} For every combination of dataset, partition and attack type, the malicious-client fraction is swept over $f/N\in\{10\%,20\%,30\%,50\%\}$. The $50\%$ row is the most demanding stress test, sitting right at the boundary that majority-based aggregation can in principle survive.

\section{Evaluation Metrics}
\label{sec:setup-metrics}

Five metrics are used in the experiments. Three of them evaluate model utility on held-out test data, while the remaining two measure how effectively the verifier distinguishes malicious clients from honest ones.

\noindent\textbf{Target accuracy (Tacc).} This metric evaluates performance only on test samples belonging to the targeted source class $c_s$:
\begin{equation}
    \mathrm{Tacc} \;=\; \frac{|\{(x, y) \in \mathcal{D}_{\text{test}} : y = c_{s},\; \hat{f}(x) = c_{s}\}|}{|\{(x, y) \in \mathcal{D}_{\text{test}} : y = c_{s}\}|}.
\end{equation}
A successful targeted attack reduces this value, whereas an effective defence keeps it close to the clean baseline. For untargeted attacks, class $0$ is taken as a fixed reference class so that results remain comparable across rows.

\noindent\textbf{Other-label accuracy (Oacc).} Other-label accuracy measures performance on all test samples whose labels are not $c_s$:
\begin{equation}
    \mathrm{Oacc} \;=\; \frac{|\{(x, y) \in \mathcal{D}_{\text{test}} : y \neq c_{s},\; \hat{f}(x) = y\}|}{|\{(x, y) \in \mathcal{D}_{\text{test}} : y \neq c_{s}\}|}.
\end{equation}
Together, $\mathrm{Tacc}$ and $\mathrm{Oacc}$ show whether the attack mainly affects the target class or also damages the rest of the task.

\noindent\textbf{Overall accuracy (ACC).} Overall accuracy is the standard test accuracy over all classes:
\begin{equation}
    \text{ACC} \;=\; \frac{|\{(x, y) \in \mathcal{D}_{\text{test}} : \hat{f}(x) = y\}|}{|\mathcal{D}_{\text{test}}|}.
\end{equation}
This is the main utility metric for untargeted attack experiments.

\noindent\textbf{True positive rate (TPR).} TPR is the fraction of malicious clients rejected by the FIP-FL audit:
\begin{equation}
    \mathrm{TPR} \;=\; \frac{|\{i : \text{client } i \text{ is malicious and rejected}\}|}{|\{i : \text{client } i \text{ is malicious}\}|}.
\end{equation}
A higher TPR means that fewer malicious clients pass through the verifier.

\noindent\textbf{False positive rate (FPR).} This is the fraction of honest clients that are wrongly rejected:
\begin{equation}
    \mathrm{FPR} \;=\; \frac{|\{i : \text{client } i \text{ is honest and rejected}\}|}{|\{i : \text{client } i \text{ is honest}\}|}.
\end{equation}
This metric captures the cost of filtering honest participants. All five metrics are averaged over three independent random seeds. Standard deviation is reported only when it is greater than $0.5$ percentage points.

\section{Implementation and Hardware}
\label{sec:setup-impl}

The implementation is written in Python~3. \texttt{NumPy} is used for linear algebra, \texttt{scikit-learn} provides the dataset loaders and feature pipelines, and Paillier encryption is implemented using the \texttt{phe} library with a $1024$-bit modulus. The FIP Verifier and Aggregation Server run as separate Python processes and communicate through in-memory queues. The clients are simulated as worker processes. Both encryption and the homomorphic inner-product computation are parallelised using \texttt{multiprocessing}.

The wall-clock results in Chapter~\ref{Comparison:chap:chap8} were measured on a Linux workstation with a $25$-core CPU and $128$~GB of RAM.
